\documentclass[10pt]{amsart}
\usepackage[a4paper,margin=23mm]{geometry}
\usepackage{amsmath,amssymb,amsthm,mathtools,hyperref,graphicx,xspace,letltxmacro,xpatch}
\usepackage[english]{babel}
\newtheorem{theo}{Theorem}[section]
\newtheorem{lemm}[theo]{Lemma}\newtheorem{prop}[theo]{Proposition}\newtheorem{coro}[theo]{Corollary}\newtheorem{conj}[theo]{Conjecture}
\theoremstyle{definition}\newtheorem{defi}[theo]{Definition}\newtheorem{rema}[theo]{Remark}\newtheorem{exam}[theo]{Example}\newtheorem{exem}[theo]{Example}\newtheorem{nota}[theo]{Notation}
\newcommand{\ie}{i.e.}\newcommand{\eg}{e.g.}\newcommand{\resp}{resp.}\newcommand{\skpt}{\par\smallskip\noindent}
\emergencystretch=3em\pagestyle{plain}
%~Mouliné par MaN_auto v.0.27.3 2023-10-24 15:01:20

\usepackage{comment}




\newcommand{\R}{\mathbb{R}}
\newcommand{\N}{\mathbb{N}}
\newcommand{\Z}{\mathbb{Z}}
\newcommand{\rmn}{\mathrm{n}}
\newcommand{\supp}{\mathrm{supp}}
\newcommand{\Ker}{\mathrm{Ker}}
\newcommand{\Tr}{\mathrm{Tr}}
\renewcommand{\Im}{\mathrm{Im}}
\newcommand{\clF}{\mathcal{F}}
\newcommand{\clA}{\mathcal{A}}
\newcommand{\di}{\displaystyle}

\newcommand{\cxi}{\langle \xi\rangle}





%%%%%%%%%%%%%%%%%%%%%%%%%%%%%%%%%%%%%%%%%%%%%%%%
\graphicspath{{./figures/}}

\newcommand*{\mk}{\mkern -1mu}
\newcommand*{\Mk}{\mkern -2mu}
\newcommand*{\mK}{\mkern 1mu}
\newcommand*{\MK}{\mkern 2mu}

%\hypersetup{urlcolor=purple, linkcolor=blue, citecolor=red}

\newcommand*{\romanenumi}{\renewcommand*{\theenumi}{\roman{enumi}}}
\newcommand*{\Romanenumi}{\renewcommand*{\theenumi}{\Roman{enumi}}}
\newcommand*{\alphenumi}{\renewcommand*{\theenumi}{\alph{enumi}}}
\newcommand*{\Alphenumi}{\renewcommand*{\theenumi}{\Alph{enumi}}}
\let\oldtilde\tilde
\renewcommand*{\tilde}[1]{\mathchoice{\widetilde{#1}}{\widetilde{#1}}{\oldtilde{#1}}{\oldtilde{#1}}}
\let\oldhat\hat
\renewcommand*{\hat}[1]{\mathchoice{\widehat{#1}}{\widehat{#1}}{\oldhat{#1}}{\oldhat{#1}}}
\let\oldforall\forall
\renewcommand*{\forall}{\mathrel{\oldforall}}
%%%%%%%%%%%%%%%%%%%%%%%%%%%%%%%%%%%%%%%%%%%%%%%%%%

\NewDocumentEnvironment{enonce}{ O{} m O{} }{\par\smallskip\refstepcounter{theo}\noindent\textbf{#2~\thetheo\if\relax\detokenize{#3}\relax\else\ (#3)\fi.}\ }{\par\smallskip}
\title{Sharp fractional Sobolev regularity for hyperbolic initial-boundary value problems}
\author{Audiard, Corentin}
\date{doi:10.5802/ahl.190; CC BY 4.0}
\begin{document}
\maketitle
\noindent\textit{Editorial note.} Only selected author result counted; complete original mathematical body and bibliography retained. Native vector figures and TikZ retained; class frontmatter/web-only definitions omitted.\par\medskip
\noindent\textit{Allowed standard prerequisites (editorial).}\begin{itemize}
\item Standard functional analysis, measure theory and the author-cited classical theorems. Standard prior results stay cited; complete source-local proofs and definitions retained.
\end{itemize}\medskip

\section{Introduction}

\subsection{Everything in a toy model}
Consider the simplest hyperbolic initial boundary value problem (IBVP)
\[
\begin{cases}
\partial_tu+\partial_xu(x,t)=0,\ (x,t)\in (\R^+)^2\\
u(x,0)=u_0(x),\\
u(0,t)=g(t)
\end{cases}
\]
When $(u_0,g)\in L^2(\R^+)^2$, the solution is piecewise defined: $u(x,t)=u_0(x-t)$ for $x-t\geq 0$, $g(t-x)$ for $x-t< 0$, it belongs to $C_tL^2$. 


It is well known that the smoothness of $(u_0,g)$ is not enough to ensure the smoothness of $u$, compatibility conditions are required: for $k\in \N$, $u\in \cap_{j=0}^kC_t^jH^{k-j}$ if and only if
\[
(u_0,g)\in \left(H^k\right)^2\text{ and }\forall\,j\leq k-1,\ u_0^{(j)}(0)=(-1)^jg^{(j)}(0).
\]
These compatibility relations are trivial here due to the solution formula, but are more generally derived considering $u$ (and its derivatives) at the corner $x=t=0$, and writing $\partial^\alpha u|_{x=0}|_{t=0}=\partial^\alpha u|_{t=0}|_{x=0}$. A basic rule of thumb being that for some regularity to hold, any compatibility condition that makes sense should be true.

For fractional regularity, not much changes except in the notoriously pathological case $s\equiv 1/2[\Z]$. Indeed even if there is no trace in $H^{1/2}(\R^+)$, the gluing of two functions in $H^{1/2}(\R^+)$ is not $H^{1/2}(\R)$. The simplest way to see this is to consider the map $f\in L^2(\R)\to f(\cdot)-f(-\cdot)\in L^2(\R^+)$. It is continuous $L^2(\R)\to L^2(\R^+)$ and $H^1(\R)\to H^1_0(\R^+)$ hence $H^{1/2}(\R)\to [L^2(\R^+),H^1_0(\R^+)]_{1/2}$ by interpolation. The interpolated space is the famous Lions--Magenes space $H^{1/2}_{00}(\R^+)$, and it is different (algebraically and topologically) from $H^{1/2}(\R^+)$: by interpolation of Hardy's inequality, any function $f\in H^{1/2}_{00}(\R^+)$ must satisfy
\[
\int_{\R^+}\frac{f^2(x)}{x}dx<\infty,
\]
this is obviously not the case for functions merely in $H^{1/2}(\R^+)$ (pick for example~$1/(1+x^2)$), see Section~\ref{notations} for more details on these spaces.


For the regularity of solutions of the BVP, this adds a ``global'' compatibility condition
\[
u\in C_tH^{1/2}(\R^+)\Leftrightarrow (u_0,g)\in H^{1/2}(\R^+)\text{ and }
\int_{\R^+}\frac{|g(x)-u_0(x)|^2}{x}dx<\infty.
\]
Our aim here is to extend these observations for general hyperbolic boundary value problems.

\subsection{Settings and results}

Let $\Omega$ be a smooth open set of $\R^d$, we consider first order boundary value problems of the form
\begin{equation}\label{IBVP}
\begin{cases}
Lu:=\left(\partial_t-\sum_{j=1}^d A_j\partial_j\right)u=0,\ (x,t)\in \Omega\times \R_t^+,\\
Bu|_{\partial\Omega}=g,\ (x,t)\in \partial\Omega\times \R_t^+,\\
u|_{t=0}=u_0,\ x\in \Omega.
\end{cases}
\end{equation}
The index $t$ in $\R_t^+$ has no meaning except to emphasize the time variable. The $A_j$'s are $q\times q$ matrices depending smoothly on $(x,t)$, $B$ is a smooth $b\times q$ matrix, $b$ is the number of boundary conditions.


For data $(u_0,g,f)\in L^2(\Omega)\times L^2(\partial\Omega\times \R_t^+)\times L^2(\R^+_t\times \Omega)$, the well-posedness of such hyperbolic BVP has been obtained in a large variety of settings, that we will only shortly mention. After the pioneering results of Friedrichs~\cite{Friedrichs} for symmetric dissipative systems, Kreiss~\cite{Kreiss} proved the well-posedness of the BVP with zero initial data in the strictly hyperbolic case ($\sum A_j\xi_j$ has only real eigenvalues of algebraic multiplicity one) under the now standard Kreiss--Lopatinskii condition on $B$. In Kreiss's framework, the case of $L^2$ initial data was then tackled by Rauch~\cite{Rauch}. Well-posedness of BVP hyperbolic with constant multiplicities was later obtained by Métivier~\cite{Met} (zero initial data), the author then proved well-posedness with $L^2$ initial data~\cite{Audiard1}. A further generalization was obtained by Métivier~\cite{Metsemigroupe} for a new class of hyperbolic operators, larger than the constant multiplicities ones. He also gave a new proof, both more general and simpler, of well-posedness with $L^2$ initial data.


For more references and results, in particular for characteristic BVP (that we do not consider here) the reader may refer to the book~\cite{Benzoni3}.


Let $\rmn$ be a normal on $\partial\Omega$, the problem~\eqref{IBVP} is said to be noncharacteristic when $\sum A_jn_j$ is invertible on $\partial\Omega$. For non characteristic boundary value problems, the main reference on the smoothness of solutions is the classical paper of Rauch and Massey~\cite{RauchMassey}, where, under no specific assumption (except of course well-posedness), the authors prove that the solution of~\eqref{IBVP} belongs to $\cap_{j=0}^k C^j_t(\R^+_t,H^{k-j}(\Omega))$ when $(u_0,g,f)\in H^k(\Omega)\times H^{k+1/2}(\partial\Omega\times\R_t^+)\times H^k(\Omega\times\R_t^+)$ and satisfy natural compatibility conditions that we describe now. For conciseness, when there is no ambiguity we will usually denote $H^k$ instead of $H^k(X)$, $X=\Omega,\partial\Omega\times \R_t^+,\Omega\times \R_t^+$.


We denote $\clA=\sum A_j\partial_j$ and define inductively $v_j$ the formal value of $(\partial_t^ju)|_{t=0}$ by
\begin{equation}\label{taylor}
\begin{split}
v_0=u_0,\ v_{j+1}&=\left(\partial_t^j\partial_tu\right)|_{t=0} =\partial_t^j(\clA u+f)|_{t=0}\\ &=\sum_{l=0}^j\binom{j}{l} \left(\partial_t^l\clA|_{t=0}\right)v_{j-l}+\partial_t^jf|_{t=0}.
\end{split}
\end{equation}
The first order compatibility condition is $Bv_0|_{\partial\Omega}=g|_{t=0}$ and the generic compatibility condition of order $j$ is
\begin{equation}\label{CCj}
\text{Compatibility at order }j:\
\partial_t^{j-1}g|_{t=0}=\sum_{l=0}^{j-1}
\binom{j-1}{l}(\partial_t^lB)v_{j-1-l}|_{\partial\Omega}.
\end{equation}
Note that~\eqref{CCj} makes sense as soon as $(u_0,g,f)\in (H^s)^3,\ s>j-1/2$. If the smoothness of the data is $j-1/2, j\in \N^*$, we define a special compatibility condition : when $\Omega=\R^{d-1}\times \R^+$, denote $x=(x',x_d)$; the condition is Compatibility at order $j-1/2$:
\begin{equation}\label{CCj-}
\partial_t^{j-1}g(x',t)-\left(\sum_{l=0}^{j-1}
\binom{j-1}{l}(\partial_t^lB)v_{j-1-l}(x',t)\right)
\in H^{1/2}_{00}\left(\R^{d-1}\times (\R^+)\right).
\end{equation}
For general smooth $\Omega$, \eqref{CCj-} is defined similarly through local maps and a partition of unity: near the boundary $\Omega$ is diffeomorphic to (a part of) $\R^{d-1}\times \R^+$ thanks to some map $\Phi$, one simply requires~\eqref{CCj-} to stand for $g(\Phi(x',0),t), (v_l\circ\Phi(x',t))_{0\,\leq\,l\,\leq\,j-1}$.


Note that due to Hardy's inequality, the $j^{\rm th}$ condition implies the condition of order $j-1/2$.
\begin{defi}
For $0\leq s<1/2$, no compatibility conditions are required, we say that any data in $(H^s)^3$ satisfy the compatibility conditions of order $s$.


If $s=k+\theta,\ -1/2< \theta <1/2,\ k\in \N^*,\
\theta\neq 1/2$, we say that data $(u_0,g,f)\in (H^s)^3$ satisfy the compatibility conditions at order $s$ when~\eqref{CCj} is satisfied for $1\leq j\leq k$.


If $s=k-1/2$, the compatibility conditions are satisfied at order $s$ when~\eqref{CCj} is true for $1\leq j\leq k-1$ and~\eqref{CCj-} is true for $j=k$.
\end{defi}

A strong $L^2$ solution of~\eqref{IBVP} is a function $u\in C_tL^2$ such that there exists a sequence $u_n$ of smooth solutions of~\eqref{IBVP} with data $(u_{0,n},g_n,f_n)$ that converge to $(u_{0},g,f)$ in $L^2$, and for any $T>0$, $\|u-u_n\| _{C([0,T],L^2)}\to 0$.

\begin{enonce}{Assumptions}
We need the smoothness of $\Omega$ and the well-posedness of~\eqref{IBVP}:
\begin{enumerate}
\item \label{assump1}
$\partial\Omega$ is a smooth hypersurface with normal $\rmn$, parametrized by local maps $(\phi_j(y'))_{1\,\leq\,j\,\leq \,J}$, $y'\in \R^{d-1}$, and $\varphi_j(y',y_d):=\phi_j(y')+y_d n(\phi_j(y'))$ are local diffeomorphisms $V_j\to U_j$, with $\varphi_j ((\R^{d-1}\times \R^{+*})\cap V_j) \subset
\Omega$, and $\cup_{j=1}^JU_j\supset \partial\Omega$. We do not assume that the $U_j$ are bounded sets, but $D\varphi_j,D\varphi_j^{-1}$ must be uniformly bounded, and $d(\Omega\setminus
\cup \Im(\varphi_j),\,
\partial\Omega)>0$.
\item \label{charac}
The boundary is uniformly not characteristic, in the sense that $\sum A_jn_j$ is invertible on $\partial\Omega$, and the inverse is uniformly bounded.
\item For data $(u_0,g,f)\in (L^2)^3$, there exists a unique strong $L^2$ solution\footnote{This assumption can be somehow weakened, as it is classical that in this framework, weak solutions are actually strong, see~\cite{LaxPhillips}.} to~\eqref{IBVP} that satisfies the semi-group estimate for $\gamma$ large enough
\begin{multline}\label{semigroupe}
\|e^{-\gamma \cdot}u\|_{C([0,t],L^2(\Omega))}+ |e^{-\gamma \cdot}u|_{\partial\Omega}|_{L^2(\partial\Omega\times [0,t])}\\
\lesssim \|u_0\|_{L^2(\Omega)}+ |e^{-\gamma \cdot }g|_{L^2(\partial\Omega\times [0,T])}
+\frac{\|e^{-\gamma \cdot}f\|_{L^2 ([0,t]\times \Omega)}}{\sqrt{\gamma}}.
\end{multline}
We use the convention that norms inside the domain are denoted $\|\cdot\|$ while norms on the boundary are denoted $|\cdot|$.
\end{enumerate}
\end{enonce}
We point out that a consequence of the semi-group estimate is the \emph{resolvent estimate}: for $\gamma$ large enough (larger than for~\eqref{semigroupe})
\begin{multline}\label{resolv}
\gamma \|e^{-\gamma t}u\|_{L^2(\Omega\times \R_t^+)}^2 +|e^{-\gamma t}u|_{\partial\Omega}|_{L^2(\partial\Omega\times \R_t^+)}^2
\\
\lesssim \left(\|u_0\|_{L^2(\Omega)}^2+ |e^{-\gamma t}g|_{L^2(\partial\Omega\times \R_t^+)}^2+\frac{
\|e^{-\gamma t}f\|_{L^2}^2}{\gamma}\right).
\end{multline}
This is readily obtained by squaring~\eqref{semigroupe} for some fixed $\gamma_0$, multiplication by $e^{-2(\gamma -\gamma_0)t},\
\gamma >\gamma_0$ and integration in $t$. Higher regularity versions of the resolvent and the semi-group estimates are a bit more delicate to state. We define weighted Sobolev spaces $H^s_\gamma$ in Section~\ref{notations}, the weighted resolvent estimate is then
\begin{equation}\label{resolvreg}
\gamma \|u\|_{H^s_\gamma}^2 +|u|_{\partial\Omega}|_{H^s_\gamma}^2
\lesssim \|u_0\|_{H^s(\Omega)}^2+ |g|_{H^s_\gamma}^2+\frac{
\|f\|_{H^s_\gamma}^2}{\gamma}.
\end{equation}
The main point of this estimate is that the $\gamma$ factor allows to absorb commutators in a priori estimates. Moreover, it implies the following (simpler to read) estimate
\begin{multline}\label{resolvreguseless}
\|e^{-\gamma t}u\|_{H^s(\Omega\times \R^+)}^2 +|e^{-\gamma t}u|_{\partial\Omega}|_{H^s(\partial\Omega\times \R^+)}^2
\\
\lesssim \|u_0\|_{H^s(\Omega)}^2+ |e^{-\gamma t}g|_{H^s(\partial\Omega\times \R^+)}^2+
\|e^{-\gamma t}f\|_{H^s(\Omega\times \R^+)}^2.
\end{multline}
We shall not need something as precise for the semi-group estimate: let $s=k+\theta,\ k\in \N,\ 0<\theta<1$, $\theta\neq 1/2$, then
\begin{multline}\label{semigreg}
\sum_{j=0}^k\left\|e^{-\gamma t}\partial_t^ju\right\|_{C\left(\R_t^+,H^{k-j+\theta}(\Omega)\right)}^2 +\left|e^{-\gamma t}u|_{\partial\Omega}\right|_{H^s\left(\Omega\times \R_t^+\right)}^2
\\
\lesssim \|u_0\|_{H^s(\Omega)}^2+ |e^{-\gamma t}g|_{H^s\left(\partial\Omega\times \R_t\right)}^2+
\left\|e^{-\gamma t}f\right\|_{H^s}^2.
\end{multline}
Both estimates should be modified when $s=k+1/2$, $k\in\N$: it is necessary to add in the right hand side the $H^{1/2}_{00}$ norm of $\partial_t^kg-\sum_0^k
\binom{k}{l}(\partial_t^lB)v_{k-1-l}$, see page~\pageref{cas1/2} for details. This is the (implicit) convention that we use in Theorem~\ref{mainth}, we refer to the proof for more details.

An interesting related feature is that the constant in $\lesssim$ can not be uniform in $\theta$, it blows up as $\theta \to 1/2$ and estimates~\eqref{resolvreg}, \eqref{semigreg} are actually not true for $\theta=1/2$.


We can now state more precisely the regularity result of Rauch and Massey:
\begin{theo}[\cite{RauchMassey}]
For any $k\in \N$, if $(u_0,g,f)\in H^k(\Omega)\times H^{k+1/2}(\partial \Omega\times \R_t^+)
\times H^k(\Omega\times \R_t^+)$ satisfy the compatibility condition up to order $k$, the solution of~\eqref{IBVP} belongs to $\cap_{j=0}^k C_t^jH^{k-j}$.
\end{theo}

The only suboptimal part of the theorem is the regularity assumption on $g$. This is due to the fact that the theorem is deduced from the homogeneous case $g=0$ with a lifting argument. It was already pointed out at the time by the authors that it could be improved (without proof), but quite unfortunately the result that remained in the literature is the suboptimal one, see for example the reference book~\cite{Benzoni3}, and in somewhat different settings the lecture notes~\cite{Met2} or the interesting discussion in the introduction of~\cite{lanigu}, where optimal results are obtained in dimension~$1$ and an integer index of regularity.


Our result is that the same property holds with boundary data in $H^k$ instead of $H^{k+1/2}$, moreover we allow $k$ to be any nonnegative real number rather than an integer.


\begin{theo}\label{mainth}
Let $s=k+\theta\in \R^+,\ k\in \N,\ 0\leq \theta<1$. If $(u_0,g,f)\in H^s(\Omega)\times H^{s}(\partial \Omega\times \R_t^+)
\times H^s(\Omega\times \R_t^+)$ satisfy the compatibility condition up to order $s$, the solution of~\eqref{IBVP} belongs for any $T>0$ to $H^s(\Omega\times [0,T])\cap (\cap_{j=0}^k C^j([0,T],H^{s-j}(\Omega)))$, and satisfies estimate~\eqref{resolvreg} and~\eqref{semigreg} for $\gamma$ large enough.
\end{theo}


The proof when $s$ is an integer is quite similar to the original argument of Rauch and Massey. Actually the fact that we handle directly nonzero boundary data leads to some slight simplifications due to the fact that it allows to avoid a reduction to the case where $B$ is constant. The fractional case is essentially an interpolation argument, however it is not trivial due to the presence of the compatibility conditions. For example, in the model case described earlier instead of interpolating $[L^2\times L^2,H^1\times H^1]_\theta$ one must identify $[L^2\times L^2,\{(u_0,g)\in H^1(\R^+)\times H^1(\R^+):\ u_0(0)=g(0)\}]_\theta$.


The literature on such problems is not very rich. Another related problem is the interpolation of Sobolev spaces with boundary conditions, that are in some sense between $H^s$ and $H^s_0$. This issue appeared quite long ago for elliptic equations on non smooth domains or parabolic problems, see e.g. the last section of~\cite{Grisvard}, \cite[Sections~14-17 of Chapter~4]{lionsmagenes2} (where most of the identification problems were left open), or the more recent (and much more involved) book~\cite{Amann2}, in particular VIII.2.5. Due to the highly technical flavour of this last reference (anisotropic Besov spaces are studied), degenerate cases (in our settings $s\in \N+1/2$) are not considered. The Schr\"odinger equation on a domain and related interpolation problems were also studied by the author in~\cite{Audiard7}, where the natural spaces for the boundary data are Bourgain spaces.

\subsection{Plan of the article}
Section~\ref{notations} is devoted to notations and a brief reminder on interpolation. The proof of Theorem~\ref{mainth} is then organized in three sections : in Section~\ref{sect3} we recall a standard smoothness result for the pure boundary value problem posed for $t\in \R$, due to Tartakov. For completeness, we include a sketch of proof that follows an argument of the (unfortunately depleted) book~\cite{ChaPi}. Theorem~\ref{mainth} in the case $s$ integer is proved in Section~\ref{sect4}. An important point is a basic lifting lemma which proves to be also useful for the general case. In Section~\ref{sect5}, smoothness is first proved for $0\leq s\leq 1$ with an interpolation argument, then for any $s$ with a non trivial differentiation argument.



\section{Notations and basic results}\label{notations}

\subsection{Basic notations}
Proofs are often reduced to the case $\Omega=\R^{d-1}\times \R^+$. In such settings, we denote the variable $x=(x',x_d)$ $x'\in \R^{d-1}$. The variables $x',t$ are called tangential, while $x_d$ is the normal variable. 


Partial differential operators acting on functions of $(x,t)$ are written as $\partial^\alpha$, $\alpha\in \N^{d+1}$, by convention $\alpha_{d+1}$ is the order of differentiation in time. A multi-index, or a differential operator, is said to be tangential when $\alpha_d=0$.


We denote $[L_1,L_2]=L_1L_2-L_2L_1$ the commutator between two linear operators.

\subsection{Sobolev spaces}
$\Omega$ is assumed to be a smooth open set as in the assumption $1$ page~\pageref{assump1}. The Sobolev spaces $H^s(\Omega)$, are defined when $s$ is an integer as
\[
\left\{u\in L^2:\ \|u\|_{H^s}^2=
\sum_{|\alpha|\,\leq\,s}\int_\Omega|\partial^\alpha u|^2dx<\infty\right\}.
\]
When $s$ is not an integer, they are defined by (complex) interpolation, $H^s=[L^2,H^k]_{s/k}$ for any integer $k$ larger than $s$. This definition does not depend on~$k$.


The Sobolev spaces for functions defined on $\Omega\times \R_t^+$ are defined in the same way, Sobolev spaces on the manifold $\partial\Omega$ (or $\partial\Omega\times [0,T]$) are defined thanks to local charts\footnote{Most references, e.g.~\cite[Chapter~1]{lionsmagenes}, assume the boundedness of $\Omega$, but everything works similarly in our settings.}.


$H^s_0(\Omega)$ is the closure of $C_c^\infty(\Omega)$ for the $H^s(\Omega)$ norm. We do have $[L^2,H^1_0]_s=H^s_0$ for $0<s<1$, except for $s=1/2$, where $H^{1/2}_{0}=H^{1/2}$ and $[L^2,H^1_0]_{1/2}=H^{1/2}_{00}$ is different algebraically and topologically from $H^{1/2}$. It is a Banach space endowed with the norm
\[
\|u\|_{H^{1/2}_{00}}^2=
\|u\|_{H^{1/2}}^2+\int_{\Omega}
\frac{|u(x)|^2}{d(x)}dx,
\]
where $d$ is the distance to $\partial\Omega$ (see~\cite{lionsmagenes}). An essential fact, regularly used in the article, is that if $X_0,X_1$ are Banach spaces, an operator $T:\ X_0\to L^2,\ X_1\to H^1_0$ maps $[X_0,X_1]_{1/2}$ to $H^{1/2}_{00}$. For example, $u\in H^s(\R^d)\to u(x',x_d)-u(x',-x_d)$ maps $H^{1/2}(\R^d)$ to $H^{1/2}_{00}(\R^{d-1}\times \R^+)$.


The weighted Sobolev spaces $H^s_\gamma$ are defined as follows :
\begin{defi}
When $s$ is a nonnegative integer we define $H^s_\gamma(\Omega\times \R_t^+)$ as the set of functions in $L^2$ such that the following norm is finite
\[
\|u\|_{H^s_\gamma}=\sum_{|\alpha|\,\leq\,s}
\left\|e^{-\gamma t}\partial^\alpha u\right\|_{L^2}.
\]
When $s$ is not an integer, $H^s_\gamma$ is defined by complex interpolation : if $k$ is an integer larger than $s$, $H^s_\gamma=[L^2_\gamma,H^k_\gamma]_{s/k}$.


$H^s_\gamma(\partial\Omega\times \R_t^+)$ is defined similarly.
\end{defi}

When $s$ is an integer, it is a straightforward consequence of Leibniz formula $\partial_t^j(e^{-\gamma t}u)=\sum \binom{j}{i} (-\gamma)^ie^{-\gamma t}\partial_t^{j-i}u$ that the $H^s_\gamma$ norm is equivalent to $\|e^{-\gamma t}u\|_{H^s}$, though with constants that depend on $\gamma$, hence the $H^s_\gamma$ spaces coincide algebraically and topologically with the set of functions such that $e^{-\gamma t}u\in H^s$.

\subsection{Traces}
Sobolev spaces on $\partial\Omega$ are defined with local maps. The trace operator is a continuous surjective morphism:
\[
\left\{
\begin{array}{ll}
H^s(\Omega)\to \prod_{k\,<\,s-1/2} H^{s-1/2-k}(\partial\Omega),\\
u\to \left(\partial_n^ku|_{\partial\Omega}\right)_{k\,<\,s-1/2},
\end{array}\right.
\]
where $\partial_n$ is the normal derivative on $\partial\Omega$.\\
For functions defined in $H^s(\Omega\times
\R^{+*})$, the trace operator on $\partial\Omega\times \R^{+*}$ and $\Omega\times \{0\}$ is more subtle, the map
\begin{equation}\label{trace2}
\left\{
\begin{array}{ll}
H^s\left(\Omega\times \R_t^{+*}\right)\to \left(\prod_{k\,<\,s-1/2} H^{s-1/2-k}\left(\partial\Omega\times \R_t^{+*}\right)\right)
\\
\mkern150mu\times \left(\prod_{k\,<\,s-1/2}H^{s-1/2-k}\left(\Omega \times \{0\}\right)\right),
\\
u\to \left(\partial_n^ku|_{\partial\Omega
\times \R^{+*}},\ \partial_t^ku|_{\Omega\times \{0\}}\right)_{k\,<\,s-1/2},
\end{array}\right.
\end{equation}
is continuous but not surjective: if $s\notin \N$,
\emph{local compatibility conditions} between $(g_k,v_k)\in (\prod H^{s-1/2-k}(\partial\Omega\times \R_t^{+*}))\times (\prod H^{s-1/2-k}(\Omega\times \{0\}))$ are required as follows (see~\cite{lionsmagenes2})
\begin{equation}\label{compaclassique1}
\forall\,k+j<s-1,\
\partial_t^jg_k|_{t=0}=\partial_n^kv_j|_{\partial\Omega}.
\end{equation}
In the case $s=1$, and $\Omega=\R^{d-1}
\times \R^{+*}$, surjectivity requires the \emph{global compatibility condition}
\begin{equation}\label{compaclassique2}
v_0(x',t)-g_0(x',t)\in H^{1/2}_{00}(\partial\Omega).
\end{equation}
This condition extends to smooth $\Omega$, see the short comment after~\eqref{CCj-}.


Provided such compatibility conditions are added, the trace map is a surjection and has a right inverse, this very well known fact will be proved later in the article in some basic cases where it is needed with more precise estimates.

\section{Regularity for the pure boundary value problem}\label{sect3}

Consider the boundary value problem
\begin{equation}\label{bvp}
\left\{
\begin{array}{ll}
Lu=f,\ (x,t)\in \Omega \times \R_t^+\\
Bu|_{\partial\Omega}=g,\\
u|_{t=0}=0.
\end{array}
\right.
\end{equation}
When $g$, $f$ can be smoothly extended by $0$ for $t<0$, the smoothness of $u$ is well known~\cite{ChaPi,Tartakoff}. The classical proof is done by first studying the pure boundary value problem posed on $t\in \R$, the case $t\in \R^+$ is then deduced by an extension by $0$ for $t<0$. We give here a minor variation of this argument that directly tackles~\eqref{bvp}.

\begin{prop}\label{regBVP}
Let $k\in \N$. If the extension of $f$ and $g$ by $0$ for $t<0$ belongs to $H^k$, then for $\gamma$ large enough the solution of~\eqref{bvp} satisfies estimate~\eqref{resolvreg} with $u_0=0,\ s=k$, in particular, $(e^{-\gamma t}u,e^{-\gamma t}u|_{\partial\Omega})\in H^k(\Omega\times \R^+_t)\times H^k(\partial\Omega\times \R_t^+)$, and $u\in H^k(\Omega\times [0,T])$ for any $0<T<\infty$.
\end{prop}


\begin{proof}
The classical plan is to straighten the boundary through local maps, then use a tangential regularization. It is done by induction on $k$, it suffices to prove the final step where we assume $u\in H^{k-1}(\R^d\times \R_t^+)$ and prove $u\in H^k$.


We fix local maps $\varphi_j$ as in Assumption~\ref{assump1}. Let $(\psi_j)_{0\,\leq\,j\,\leq\,J}$ be a partition of unity associated to $\Omega\cup (\cup_j\Im(\varphi_j))$.

%\begin{comment}
%It is sufficient to prove that for any $j$, $\psi_je^{-\gamma t}u
%\in H^k$, ou par les règles de composition usuelles $\Phi_j^*(\psi_je^{-\gamma t}u)\in H^k$, où $^*$ est le pullback. En multipliant l'équation par $e^{-\gamma t}$ on se ramène aux Sobolev usuels, \\
%L'équation transportée est
%\[
%\psi_j\left(\partial_tu_j+\sum_i\left(\sum_k A_k(\Phi_j(y))(D_y\Phi_j)^{-1}_{ik}(y)\right)
%\partial_{y_i}\right)u(\Phi_j(y))=\psi_jf(\Phi_j(y)).
%\]
%\end{comment}

\noindent We denote the new variable $y=(y',y_d)$, $u_j=(\psi_je^{-\gamma t}u)
\circ\varphi_j$, and $u_0=\psi_0 u$, $L_j=\partial_t+\gamma+\sum_i\left(\sum_k A_k(D_y\varphi_j)^{-1}_{ik}(y)\right)
\partial_{y_i}$. For $1\leq j\leq J$, $u_j$ satisfies
\begin{equation}\label{eqredresse}
\left\{
\begin{array}{ll}
L_ju_j+\left([\psi_j,L]e^{-\gamma t}u\right)\circ \varphi_j=e^{-\gamma t}(\psi_jf)\circ \varphi_j\\
:=f_j, \left(y',y_d,t\right)\in \R^{d-1}\times \R^+\times \R_t^+,
\\
B(\varphi_j(y',0))u_j(y',0,t)=e^{-\gamma t}(\psi_j g)(\varphi_j(y',0),t):=g_j.
\end{array}
\right.
\end{equation}
This is still a non characteristic BVP, indeed $D\varphi_j(y',0)=
(\begin{smallmatrix} D_{y'}\phi_j & \rmn
\end{smallmatrix})$, the
last column $\rmn$ is a vector orthogonal to the previous columns hence $(D\varphi)^{-1}$ is of the form $
(\begin{smallmatrix} A\\ \rmn^t
\end{smallmatrix})$, in particular the matrix in factor of
$\partial_{y_d}$ is $\sum_kA_k n_k$, which is invertible by assumption.


For simplicity we still denote $B$ for $B\circ \varphi_j(\cdot,0)$.

%\begin{comment}
%Note that the boundary $y_d=0$ is not characteristic for $L_j$, indeed the last line of $(D\varphi_j)^{-1}$ is parallel to the normal, say $\lambda \rmn$, $\lambda\neq 0$. and the coefficient of $\partial_{y_d}u_j$ is
%\[
%\sum_{k=1}^q A_k(D_y\Phi)^{-1}_{dk}=\lambda
%\sum A_k n_k\text{, invertible}.
%\]
%\end{commment} 

The regularization procedure was introduced by H\"ormander~\cite{hormanderBVP}: for $v\in L^2(\R^p),\ p\geq 1$, define
\[
\|v\|_{H^{s,\delta}(\R^p)}^2=\int_{\R^p}|\widehat{v}(\xi)|^2
\frac{\left(1+|\xi|^2\right)^{s+1}}{1+|\delta \xi|^2}d\xi
\longrightarrow_{\delta\,\to\,0} \|v\|_{H^{s+1}}^2.
\]
Let $\rho\in C_c^\infty(\R^{p})$, such that $|\widehat{\rho}(\xi)|\lesssim |\xi|^{m},
\ m>k$ and $\widehat{\rho}$ does not cancel on a neighborhood outside $0$ (for example take $m>k$ even, $\rho_0\in C_c^\infty$ with $\int_{\R^p}\rho_0\,dx\neq 0$ and set $\rho=\Delta^{m/2}(\rho_0)$). Define $\rho_\varepsilon=\rho(\cdot/\varepsilon)/\varepsilon^p$. It is an exercise in calculus that for $0\leq s\leq k-1$, an equivalent norm to $\|\cdot\|_{H^{s,\delta}}$ -\emph{uniformly in} $\delta$- is
\begin{equation}\label{equivsobo}
\|v\|_{L^2}+\left(\int_0^1\|v*\rho_\varepsilon\|_{L^2}^2\frac{1} {\varepsilon^{2(s+1)}(1+\delta^2/\varepsilon^2)}\frac{d\varepsilon}{\varepsilon}\right)^{1/2}\sim
\|v\|_{H^{s,\delta}(\R^p)}.
\end{equation}
Friedrich's lemma can be generalized in such settings: for $P$ a first order differential operator with smooth coefficients
\begin{equation}\label{friedrichs}
\int_0^1\|[P,\rho_\varepsilon*]v\|_{L^2}^2\frac{1} {\varepsilon^{2(s+1)}(1+\delta^2/\varepsilon^2)}\frac{d\varepsilon}{\varepsilon}\lesssim \|v\|_{H^{s,\delta}}^2.
\end{equation}
For details, we refer to~\cite[Chapter~2, Section~6]{ChaPi}.


We shall use tangential mollifiers $\rho_\varepsilon(x',t)$ for the functions $u_j$, $1\leq j\leq J$, and full mollifiers $\rho_\varepsilon(x,t)$ for $u_0$ in the following way : $u_0$ is extended by $0$ outside $\Omega$, for $\varepsilon$ small enough, $\rho_\varepsilon*u_0$ is supported in $\Omega\times\R_t$, we choose such $\varepsilon$.


Everything in~\eqref{eqredresse} is extended by $0$ for $t<0$. Note that due to the assumptions on $f,g$, the extensions of $(f_j,g_j)$ are still in $H^k$. We apply $\rho_\varepsilon*$ to~\eqref{eqredresse} for $1\leq j\leq J$:
\begin{equation}\label{eqreg}
\left\{
\begin{array}{ll}
L_j\rho_\varepsilon*u_j=\rho_\varepsilon*f_j-\rho_\varepsilon*\left[\psi_j,L_j\right]e^{-\gamma t}u\circ
\varphi_j-\left[\rho_\varepsilon*,L_j\right]e^{-\gamma t}u_j,\\
B\left(\rho_\varepsilon*u_j\right)|_{y_d=0}=\rho_\varepsilon*g_j -\left[\rho_\varepsilon*,B]u_j\right|_{y_d=0}.
\end{array}
\right.
\end{equation}
We have $\rho_\varepsilon*u_j\in L^2(\R^+_{y_d},H^\infty(\R^{d-1}\times \R_t^+))$, its $H^1$ regularity is deduced from the boundary being non characteristic with the following standard argument: the right hand side of the first equation in~\eqref{eqreg} is in $L^2(\R^{d-1}\times
\R^+\times \R_t^+)$, hence
\begin{equation}\label{characgain}
\partial_{y_d}(\rho_\varepsilon u_j)=T\left(\rho_\varepsilon* u_j\right)+R_j,
\end{equation}
with $T$ a tangential differential operator and $R_j\in L^2$. This implies $\rho_\varepsilon*u_j
\in H^1$, and we can use the resolvent estimate~\eqref{resolv}:

%\begin{comment}
%Clairement pour $\delta>0$, cette quantité est finie ssi $u\in H^{k-1}$, et $\|u\|_{H^k}=\lim_{\delta\to 0}
%\|u\|_{k-1,\delta}$. On peut vérifier par calcul que
%\[
%\|u\|_{k-1,\delta}^2\sim \|u\|_{k-1}^2+\int_0^1|u*\rho_\varepsilon|^2\varepsilon^{-2k}(1+\delta^2/\varepsilon^2)^{-1}\frac{d\varepsilon}{\varepsilon}.
%\]
%\end{comment}

\begin{multline*}
\gamma \left\|\rho_\varepsilon*u_j\right\|_{L^2}^2 +\left|\rho_\varepsilon*u_j\right|_{L^2}^2 \\
\begin{aligned}
&\lesssim
\frac{\left|\rho_\varepsilon*f_j\right\|_{L^2}^2 +\left\|\rho_\varepsilon*\left[\psi_j,L_j\right]e^{-\gamma t}u\circ\varphi_j\right\|_{L^2}^2+\left\|\left[\rho_\varepsilon*,L_j \right]u_j\right\|^2_{L^2}}{\gamma}\\
& +\left|\rho_\varepsilon*g_j-\left[\rho_\varepsilon*,B\right]u_j\right|_{L^2}^2.
\end{aligned}
\end{multline*}
Multiplying by $\varepsilon^{-2k-1}
\left(1+(\delta/\varepsilon)^2
\right)^{-1}$, integrating in $\varepsilon$ and using Friedrich's lemma~(inequality~\eqref{friedrichs}) with $s=k-1$, we have
\begin{equation}\label{estimtan}
\gamma \|u_j\|_{L^2H^{k-1,\delta}}^2+ |u_j|_{H^{k-1,\delta}}^2\lesssim
\frac{\|f_j\|_{H^k}^2+\left\|\left[\psi_j,L_j\right] e^{-\gamma t}u\circ \varphi_j\right\|_{L^2H^{k-1,\delta}}^2}{\gamma}+ \|g_j\|_{H^k}^2.
\end{equation}
The commutator $[\psi_j,L_j]$ is the multiplication by a smooth matrix $\theta_j$, note that since $\psi_j$ only depends of the space variables, $[\psi_j,\partial_t]=[\psi_j,\gamma]=0$, in particular is independent of $\gamma$. Due to the special structure of the local maps, $\varphi_i^{-1}
\circ\varphi_j$ has the form $(\varphi_{i,j}(y'),y_d)$ hence
\[
\theta_j e^{-\gamma t}u\circ \varphi_j=
\sum_1^J \psi_i\theta_ju_i(\varphi_{i,j}(y'),y_d)+\theta_ju_0\circ \varphi_j.
\]
Thanks to composition rules (in $H^{s,\delta}$, again see~\cite{ChaPi}),
\[
\left\|\left[\psi_j,L_j\right]e^{-\gamma t}u\circ \varphi_j\right\|_{L^2H^{k-1,\delta}}\lesssim \sum_{i=1}^J \|u_i\|_{L^2H^{k-1,\delta}} +\|u_0\|_{H^{k-1,\delta}}
\]
For $\gamma$ large enough, this can be absorbed in (the sum over $j$ of) the left-hand side of~\eqref{estimtan}:
\begin{equation}\label{estimtan2}
\sum_{j=1}^J\gamma \|u_j\|_{L^2H^{k-1,\delta}}^2+ |u_j|_{H^{k-1,\delta}}^2\lesssim
\frac{\sum_1^J\|f_j\|_{H^k}^2+\|u_0\|_{H^{k-1,\delta}}^2}{\gamma}+ \sum_1^J|g_j|_{H^k}^2.
\end{equation}
It seems ``moral'' that noncharacteristicity should imply the same bound for $\|u_j\|_{H^{k-1,\delta}}$, however the $H^{k-1,\delta}$ norm is a non local norm for functions defined on $\R^d\times \R_t$, hence such an assertion is not clear. Instead we first obtain interior estimates with similar, simpler computations
\begin{multline}\label{interieur}
\gamma \|u_0\|_{H^{k-1,\delta}}^2
\\
\lesssim \frac{\|f_0\|_{H^k}^2 +\left\|e^{-\gamma t}\widetilde{\psi_0}u\right\|_{H^{k-1,\delta}}^2}{\gamma},\ \supp(\widetilde{\psi_0)}\subset \Omega,\
\widetilde{\psi_0}\equiv 1\text{ on }
\supp(\psi_0).
\end{multline}
Decomposing again $\widetilde{\psi_0}u =\sum_{j=0}^J \widetilde{\psi_0}\psi_ju$, and following the same lines that led to~\eqref{estimtan2},
\begin{equation}
\sum_{j=1}^J\gamma \left\|\widetilde{\psi_0}
\psi_j u\circ \varphi_j\right\|_{L^2H^{k-1,\delta}}^2 \lesssim
\frac{\sum_1^J\|f_j\|_{H^k}^2+\|u_0\|_{H^{k-1,\delta}}^2}{\gamma}+\sum_1^J|g_j|_{H^k}^2.
\end{equation}
A simple consequence of the definition of the $H^{s,\delta}(\R_{y_d}\times \R^{d-1}\times \R_t)$ spaces is that for any tangential differential operator $D$ of order $1$ and $s\geq 1$
\begin{equation}\label{transfert}
\|Dv\|_{H^{s-2,\delta}}\lesssim
\frac{1}{C}\|v\|_{H^{s-1,\delta}}+C\|v\|_{L^2\left(\R_{y_d},H^{s-1,\delta}\left(\R_t\times \R^{d-1}\right)\right)}.
\end{equation}
Now for $j\geq 1$, each function $\widetilde{\psi_0}\psi_ju\circ \varphi_j$ is compactly supported in $\R^{d-1}\times \R^{+*}\times \R_t$, its extension by zero to $y_d<0$ is smooth, and on its support $L_j$ is (uniformly) non characteristic, hence we can use an argument similar to the one for the $H^1$ regularity for $\rho_\varepsilon*u_j$ (see~\eqref{characgain}) combined with estimate~\eqref{transfert} to deduce
\begin{equation}\label{interieur2}
\sum_{j=1}^J\gamma \left\|\widetilde{\psi_0}
\psi_j u\circ \varphi_j\right\|_{H^{k-1,\delta}}^2
\lesssim
\frac{\|f_j\|_{H^k}^2+\|u_0\|_{H^{k-1,\delta}}^2}{\gamma} +\sum_1^J|g_j|_{H^k}^2 +\gamma \|u\|_{H^{k-1}_\gamma}^2.
\end{equation}
Note that the term $\gamma\|e^{-\gamma t}u\|_{H^{k-1}_\gamma}^2$ is present due to the factor $\gamma$ in the definition of $L_j$. Thanks to the induction assumption, this lower order term is bounded by $\|g\|_{H^{k-1}_\gamma}^2+\|f\|_{H^{k-1}_\gamma}^2$. Putting together~\eqref{estimtan2}, \eqref{interieur}, \eqref{interieur2} we have
\[
\left(\sum_1^J\|u_j\|_{L^2H^{k-1,\delta}}^2+\|u_0\|_{H^{k-1,\delta}}^2\right) +\sum_1^J|u_j|_{H^{k-1,\delta}}^2
\lesssim \left\|e^{-\gamma t}f\right\|_{H^k}^2 +|e^{-\gamma t}g|_{H^k}^2.
\]
Letting $\delta\to 0$ we have $u_j\in L^2H^k,\ 1\leq j\leq J$ and $u_0\in H^k$. Normal regularity is finally gained thanks to the uniform non characteristicity. Estimate~\eqref{resolvreg} is then easily obtained by differentiation (which can now be done) and use of the $L^2$ resolvent estimate~\eqref{resolv}.
\end{proof}

\section{Smoothness of the IBVP: the integer case}\label{sect4}

We assume in this section that $(u_0,g,f)\in (H^k)^3$ satisfy the compatibility conditions~\eqref{CCj} up to order $k$, and we prove Theorem~\ref{mainth} in these settings.


To prove that $u\in \cap_{j=0}^k C_t^jH^{k-j}$, the strategy is to use the regularity for the pure boundary value problem by subtracting an approximate solution (actually a Taylor expansion at $t=0$) to $u$. For technical reasons, it is necessary to use much more regular data that satisfy compatibility conditions to higher order. The construction of such data requires the following lifting lemma that is also used in the next section.

\begin{lemm}\label{releve}
For $m\in \N$, there exists a lifting map $R_m:\ H^{s}(\partial\Omega)\to H^{m+s+1/2}\linebreak(\partial\Omega\times \R_t),$ continuous for any $s>0$ such that
\begin{equation}\label{trace}
\partial_t^{m}R_mg|_{t=0}=g,\ \partial_t^jR_mg|_{t=0}=0,\ j<m+s,
\ j\neq m,
\end{equation}
and for $r<m+1/2$, $|\|R_m|\|_{L^2\to H^r}<<1$ is arbitrarily small.
\end{lemm}

\begin{proof}
Up to the use of local maps, the problem is reduced to $\partial\Omega=\R^{d-1}$, and to construct a lifting valued in $H^{m+s+1/2}(\R^{d-1}\times \R_t)$. The variables are denoted $(x',t)$.


We choose $\chi\in C_c^\infty(\R)$ such that $\chi^{(k)}(0)=0,\ k\neq m,\ \chi^{(m)}(0)=1$. We use the Fourier transform on $\R^{d-1}\times \R_t$ and denote $\xi$ the dual variable of $x'$, $\tau$ the dual variable of $t$, and $\lambda$ is a large parameter to fix later:

\begin{multline*}
\widehat{R_mg}=\frac{\widehat{\chi}(\tau/(\lambda\langle \xi\rangle))}{(\lambda \cxi)^{m+1}}\widehat{g}(\xi),\\
\text{ equivalently }
\clF_{x'}\left(R_m(g) (\xi,t)\right)=\frac{\chi(\lambda t\langle \xi\rangle)}{\lambda^{m}\langle \xi\rangle^{m}}\widehat{g}(\xi),\
\cxi=\sqrt{1+|\xi|^2}.
\end{multline*}
The trace relations~\eqref{trace} are obvious from the second formula. The $H^{m+s+1/2}$ norm is easily bounded
\begin{align*}
\|R_mg\|_{H^{m+s+1/2}(\R^d)}^2&=\int \frac{\left|\widehat{\chi}(\tau/(\lambda\langle \xi\rangle))\right|^2|\widehat{g} (\xi)|^2} {(\lambda \cxi)^{2(m+1)}}\left(\cxi^2+\tau^2\right)^{m+s+1/2}d\xi d\tau
\\
&=\int \frac{\left|\widehat{\chi}(\tau)\right|^2|\widehat{g}(\xi)|^2}{(\lambda \cxi)^{2(m+1)}} \left(\cxi^2\left(1+\lambda^2\tau^2\right)\right)^{m+s+1/2} d\xi \lambda \cxi d\tau
\\
&\leq \int \left|\widehat{g}(\xi)\right|^2\cxi^{2s}
\int\left|\widehat{\chi}(\tau)\right|^2
\frac{\left(1+\lambda^2\tau^2\right)^{m+s+1/2}}{\lambda^{2m+1}}d\tau\,d\xi.
\\
&\lesssim \lambda^{2s}\|g\|_{H^{s}}^2.
\end{align*}
With the same computation
\[
\|R_mg\|_{H^r}^2\leq
\int \frac{|\widehat{g}|^2}{\cxi^{2(m-r)+1}}\int \frac{|\widehat{\chi}(\tau)|^2\left(1+\lambda^2
\tau^2\right)^r}{\lambda^{2m+1}}d\tau\,d\xi\lesssim
\frac{\|g\|_{L^2}^2}{\lambda^{2(m-r)+1}}.
\]
It is therefore sufficient to choose $\lambda$ large enough to ensure the smallness of\\$\|R_m\|_{L^2\,\to\,H^r}$.
\end{proof}

\begin{lemm}[Construction of smooth compatible data]\label{compahaute}
Let $k\geq 0$, $(u_0,g,f)\in (H^k)^3$ satisfying the compatibility conditions up to order $k$. For any $m>k$, there exists $(u_{0,n},g_n,f_n)\in (H^\infty)^3$ satisfying the compatibility conditions up to order $m$, and such that
\[
\left\|(u_0,g,f)-(u_{0,n},g_n,f_n)\right\|_{(H^k)^3}
\to 0.
\]
\end{lemm}


\begin{proof}
By density of smooth functions, there exists a sequence $(u_{0,n},g_n,f_n)\in (H^\infty)^3$ converging to $(u_0,g,f)$ in $(H^k)^3$. We denote $v_{j,n}$ the corresponding functions in~\eqref{taylor}. For $j\geq 1$ the ``compatibility error'' is defined as
\[
\varepsilon_{j,n}:=
\partial_t^{j-1}g_n|_{t=0}
-\sum_{l=0}^{j-1}\binom{j}{l} \left(\partial_t^lB\right)v_{j-1-l,n}|_{\partial\Omega}.
\]
Due to the compatibility conditions and continuity of traces we have
\[
\forall\,1\leq j\leq k,\
\left\|\varepsilon_{j,n}\right\|_{H^{k-j+1/2}}\longrightarrow_{n} 0.
\]
As a consequence, given a lifting operator $R_{j-1}$ as in Lemma~\ref{releve}, $\|R_{j-1}\varepsilon_{j,n}\|_{H^k}\to_n 0$.\\
For $k<j\leq m$, $\varepsilon_{j,n}$ is not small in any Sobolev space, nevertheless from Lemma~\ref{releve} there exists a lifting $R_{j-1,n}$ such that $\|R_{j-1,n}\varepsilon_{j,n}\|_{H^k}
\leq 1/n$. We then define
\[
\widetilde{g_n}:=g_n-\sum_{j=1}^m
R_{j-1,n}(\varepsilon_{j,n}).
\]
This choice ensures that compatibility conditions are satisfied by $(u_{0,n},\widetilde{g_n},f_n)$ up to order $m$ and $\|\widetilde{g_n}-g\|_{H^k}
\to 0$.
\end{proof}

\subsection{Proof of Theorem~\ref{mainth} (integer case)}
\begin{proof}
We follow the notations of Lemma~\ref{compahaute}; $v_{j,n}$ are smooth functions defined by~\eqref{taylor} for smooth data $(u_{0,n},g_n,f_n)$. We define the approximate solution
\[
u_{app,n}(x,t)=\sum_{j=0}^{m-1}\frac{t^j}{j!} v_{j,n}(x)\chi(t),\ \chi\in C_c^\infty (\R^+),\ \chi\equiv 1\text{ near 0}.
\]
We solve then
\[
\left\{
\begin{array}{ll}
Lw_n=f_n-Lu_{app,n},\\
w_n|_{t=0}=0,\\
Bw_n=g_n-Bu_{app,n},
\end{array}
\right.
\]
By construction, the data $(0,g_n-Bu_{app,n},f_n-Lu_{app,n})$ are smooth and it is easily seen that $\partial_t^j(g_n-Bu_{app,n})=0,$ $\partial_t^j(f_n-Lu_{app,n})=0,\ j\leq k+1$ provided $m\geq k+4$. Hence according to Proposition~\ref{regBVP}, the solution $w_n$ belongs to $H^{k+2}$, this implies by Sobolev embedding $w_n\in\cap_{j=0}^{k+1}C_t^jH^{k+1-j}$. Therefore $u_n:=w_n+u_{app,n}$ is also in $\cap_{j=0}^{k+1}C_t^jH^{k+1-j}$, and it is a solution of~\eqref{IBVP} with data $(u_{0,n},g_n,f_n)$.


Using a differentiation argument similar to the proof of Proposition~\ref{regBVP}, but much simpler since no regularization is needed, we see that $u_n$ satisfies~\eqref{semigreg}:
\begin{multline*}
\sum_{j=0}^k \left\|\partial_t^j\left(e^{-\gamma t}u_n\right)\right\|_{C\left(\R^+,H^{k-j}(\Omega)\right)}+ \left|e^{-\gamma t}u_n\right|_{\partial\Omega}|_{H^k}\\
\lesssim \left(\|u_{0,n}\|_{H^k(\Omega)}+ \left|e^{-\gamma t}g_n\right|_{H^k(\partial\Omega\times \R_t}
+\left\|e^{-\gamma t}f_n\right\|_{H^k}\right),
\end{multline*}
as well as~\eqref{resolvreg}. The same estimates, applied to $u_p-u_q,\ (p,q)\in \N^2$, show that $(u_n)$ is a Cauchy sequence in $\cap_{j=0}^kC_t^jH^{k-j}$, but since $(u_n)$ converges (in $L^2$) to the solution $u$ of~\eqref{IBVP} with data $(u_0,g,f)$, this ensures that $u\in
\cap_{j=0}^kC_t^jH^{k-j}$. The estimate~\eqref{resolvreg} is then an elementary differentiation argument : tangential regularity is obtained directly by differentiation (which is now legal) and use of the $L^2$ estimate, while normal regularity uses the non characteristicity.

%\begin{comment}
%Existence d'une solution approchée : partons de $\partial_tu-Au=f$, par calcul direct on a $\partial_t^ju=B_ju+C_jf$, $B_j$ opérateur différentiel spatial d'ordre $j$, $C_j$ d'ordre $j-1$. On pose
%\[
%u_{app}=\sum_{j=0}^N(B_ju_{0,n}+C_jf_n|_{t=0})
%\frac{t^j}{j!}\chi(t),\ \chi\in C_c^\infty,
%\chi\equiv 1\text{ près de }0.
%\]
%On a $u_{app}\in H^\infty(\Omega\times \R_t)$,
%\[
%\left\{
%\begin{array}{ll}
%L(u_n-u_{app})=f_n-Lu_{app},\\
%u_n-u_{app}|_{t=0}=0,\\
%B(u_n-u_{app}=g_n-Bu_{app},
%\end{array}
%\right.
%\]
%Par calcul, on voit que $\partial_t^j(g_n-Bu_{app})=0,$ $\partial_t^j(f_n-Lu_{app})=0,\ j\leq N$ donc $(g_n-Bu_{app},f_n-Lu_{app})\in H^{N+1}_0$. Par régularité pour le problème aux limites pur (Proposition~\ref{regBVP}) on a donc $u_n-u_{app}
%\in H^{N+1}(\R^+\times \Omega)$ $(u_n-u_{app}$ est solution régulière par unicité de la solution $L^2$).
%
%En conséquence, $u_n\in H^{N+1}$, et l'estimée a priori $L^2$ est disponible. Par les manipulations habituelles de dérivation tangentielle et noncaractéristicité du bord, on obtient l'estimée de semi groupe $H^N$
%\[
%\sum_{j=0}^N \|\partial_t^ju_n\|_{L^\infty_TH^{N-j}}
%\lesssim \|u_{0,n}\|_{H^N}+\|g_n\|_{H^N}+\|f_n\|_{H^N}.
%\]
%Pour $N=k$, l'estimation appliquée à $u_p-u_q$ montre que $u_n$ est une suite de Cauchy $L^\infty_TH^k$, et donc $u\in L^\infty_TH^k$ (unicité de la limite). La preuve est déjà terminée. Apparemment $N=k$ suffisait.
%\end{comment}

\end{proof}

\section{Regularity for positive $s$}\label{sect5}
For ease of presentation, we only detail the case $\Omega =\R^{d-1}\times \R^+$. The general case can be obtained by using a partition of unity as in the previous section. 


In this section, we follow the (non standard) convention that $H^s_0$ is $H^{1/2}_{00}$ if $s=1/2$.

Under such settings, we can assume that $A_d$ is invertible and $A_d^{-1}$ is uniformly bounded. Furthermore since $B:\ \R^q\to
\R^b$ has maximal rank $b$, there exists a smooth basis of $\Ker\,B$ (as a smooth vector bundle over the contractible space $\R^{d-1}\times \R_t^+$) that we denote $(k_1,\cdots k_{q-b})$. A basis $(v_j)_{1\,\leq\,j\,\leq\,b}$ of $(Ker B)^\perp$ is then obtained easily:
\[
\widetilde{B}=
\begin{pmatrix}
B\\
k_1^t\\
\vdots\\
k_{q-b}^t
\end{pmatrix}\text{ is an isomorphism }
\R^q\to \R^q,\text{ we can choose }
v_j=\widetilde{B}^{-1}(e_j),\ 1\leq j\leq b.
\]
We remind that compatibility conditions of order $s=k+\theta,\ k\in \N^*,\ 0<\theta<1$ are defined as follows:
\begin{enumerate}
\item If $\theta <1/2$, then compatibility conditions~\eqref{CCj} up to order $k$ are satisfied.
\item If $\theta>1/2$, then compatibility conditions~\eqref{CCj} up to order $k+1$ are satisfied.
\item If $\theta=1/2$, compatibility conditions up to order $k$ are satisfied and
\begin{multline*}
\int_{\R^{d-1}} \left|\partial_t^{k-1} g(x',x_d)-\sum_{j=0}^{k-1} \binom{k-1}{j}\left(\partial_t^jB\right)\left(A_{k-1-j}u_0+B_{k-1-j}f|_{t=0}\right)|(x',x_d)\right|^2\\
\frac{dx_d}{x_d}<\infty.
\end{multline*}
\end{enumerate}


\subsection{The case $0<s<1$} 
We define
\[
\left\{
\begin{array}{ll}
X_0=L^2(\Omega)\times L^2\left(\partial\Omega\times \R_t^+\right),\\
X_1=\left\{(u_0,g)\in H^1(\Omega)\times H^1\left(\partial\Omega\times \R_t^+\right):
\ Bu_0|_{\partial\Omega}=g|_{t=0}\right\}\times H^1
\end{array}
\right.
\]
From the previous section, the map $(u_0,g,f)\to u$ solution of~\eqref{IBVP} is continuous
\begin{align*}
X_0\times L^2&\to C_tL^2\text{ and }\\
X_1\times H^1\left(\Omega\times \R_t^+\right)&\to C_tH^1\cap C^1_tL^2.
\end{align*}
Let us define for $0\leq\theta \leq 1$
\[
X_\theta=\left\{(u_0,g)\in \left(H^\theta\right)^2:\ \text{the compatibility condition of order }\theta\text{ is satisfied}\right\},
\]
(note that compatibility conditions of order less than $3/2$ do not involve $f$).


Both the semi-group estimate~\eqref{semigreg} and the resolvent estimate~\eqref{resolvreg} follow from an interpolation argument if we can prove that
\begin{equation}\label{interpX}
X_\theta=[X_0,X_1]_\theta.
\end{equation}
More precisely, since the resolvent estimate implies for $s=0,1$
\begin{align*}
\gamma\|u\|_{L^2_\gamma}^2+\|u|_{\partial\Omega}\|^2_{L^2_\gamma}\lesssim \left\|\left(u_0,e^{-\gamma t}g\right)\right\|_{X_0}^2 &+\frac{\|f\|_{L^2_\gamma}^2}{\gamma}\\
\gamma\|u\|_{H^1_\gamma}^2+\|u|_{\partial\Omega}\|^2_{H^1_\gamma}\lesssim C(\gamma)\left\|\left(u_0,e^{-\gamma t}g\right)\right\|_{X_1}^2 &+\frac{\|f\|_{H^1_\gamma}^2}{\gamma},
\end{align*}
the interpolation identity~\eqref{interpX} implies
\begin{equation}\label{resolvsupersharp}
\gamma\|u\|_{H^\theta_\gamma}^2+\|u|_{\partial\Omega}\|^2_{H^\theta_\gamma}\lesssim C'(\gamma)\left\|\left(u_0,e^{-\gamma t}g\right)\right\|_{X_\theta}^2 +\frac{\|f\|_{H^\theta_\gamma}^2}{\gamma}.
\end{equation}
(a better estimate would require to use weighted $X^\theta$ spaces, a course that we chose not to follow).


\subsection{Proof of~\eqref{interpX}}
\begin{proof}
Consider the map $(x\in \Omega\to u_0(x))\to (x\in \Omega\to
\widetilde{B}(x',x_d)u_0(x)\linebreak :=\widetilde{u_0}(x))$. Since $\widetilde{B}$ is smooth and invertible for any $(x',x_d)\in \R^{d-1}\times \R^+$, the map is an isomorphism $(H^s(\Omega))^q\to (H^s(\Omega))^q$, and the compatibility condition of first order can be rewritten
\[
Bu_0|_{\partial\Omega}=g|_{t=0}
\Leftrightarrow B\widetilde{B}^{-1}
\widetilde{B}u_0|_{\partial\Omega} =g|_{t=0}\Leftrightarrow
\begin{pmatrix}
I_b & 0
\end{pmatrix}
\widetilde{u_0}|_{\partial\Omega} =g|_{t=0},\text{ with }\widetilde{u_0}=\widetilde{B}u_0,
\]
$I_b$ is the identity matrix of size $b$. Similarly the compatibility condition of order $1/2$ is $
(\begin{smallmatrix}
I_b & 0
\end{smallmatrix})\widetilde{u_0}(x)-g(x',x_d)\in H^{1/2}_{00}(\R^{d-1}\times \R^+)$\\
This transformation ``diagonalizes''~\eqref{interpX} in the following sense : set for $0\leq \theta\leq 1$
\[
Y_\theta=
\left\{
\begin{array}{ll}
\di \left\{(u_0,g)\in H^\theta\times H^\theta:\ u_0|_{x_d=0}=g|_{t=0}\right\}\text{ if }\theta >1/2,
\\
\di \left\{(u_0,g)\in H^\theta\times H^\theta:\
\int_{\R^{d-1}\times \R^+}\frac{\left|u_0(x',x_d)-g(x',x_d)\right|^2}{x_d}dx<\infty\right\}\text{ if }\theta =1/2,\\
\di H^\theta(\Omega)\times H^\theta(\partial\Omega\times \R_t),\ \text{ if }\theta <1/2,
\end{array}\right.
\]
where $u_0$ and $g$ are now \emph{scalar functions}. We are reduced to prove
\[
\left[L^2\times L^2,H^1\times H^1\right]_\theta=H^\theta\times H^\theta,\ [Y_0,Y_1]_\theta=Y_\theta.
\]
By definition $[L^2,H^1]_\theta=H^\theta$, so the first equality is trivial. In the second case, surprisingly, we were not able to find results in the literature except in the simplest case $\theta<1/2$, which is in~\cite[Section~14]{lionsmagenes2}.
\end{proof}

\begin{lemm}\label{interpfacile}
For $\theta<1/2$, $[Y_0,Y_1]_\theta=Y_\theta$.
\end{lemm}


\begin{proof}
The following inclusions are clear : $H^1_0\times H^1_0\subset Y_1\subset H^1(\Omega)
\times H^1(\partial\Omega\times \R^+)$. On the other hand, for $\theta<1/2$ we have $[L^2,H^1_0]_\theta=H^\theta$ (\cite[Chapter~1 Section~11]{lionsmagenes}), and we can conclude
\[
H^\theta\times H^\theta=\left[L^2\times L^2,H^1_0\times H^1_0\right]_\theta\subset \left[Y_0,Y_1\right]_\theta\subset \left[L^2\times L^2,H^1\times H^1\right]_\theta= H^\theta\times H^\theta.
\]
\end{proof}

\begin{lemm}\label{relevecoin}
For $0<\theta\leq 1$, there exists a universal (independent of $\theta$) operator~$R$
\[
R:\ Y_\theta\to H^{\theta+1/2}(\Omega\times \R^+),\ \forall\, 0<\theta\leq 1.
\]
\end{lemm}


\begin{proof}
This is a result due to Grisvard~\cite{Grisvard}, for completeness we include a simple proof. Given $(u_0,g)\in (H^\theta)^2$, from Lemma~\ref{releve} there exists an operator $R_b:\ g\to R_b(g)\in H^{\theta+1/2}$ which is independent of $\theta$. By construction, $R_bg|_{t=0}-u_0\in H^\theta_0$. If $\theta=1/2$, we also notice
\[
R_bg(x',y,0)-u_0(x',x_d)=
\underbrace{R_bg(x',x_d,0)-g(x',x_d)}_{H^{1/2}_{00}\text{ by interpolation}} +\underbrace{g(x',x_d)-u_0(x',x_d)}_{H^{1/2}_{00}
\text{by assumption}}.
\]
If there exists an universal lifting $R_0:\ H^\theta_0(\Omega)\to \{u\in H^{\theta+1/2}(\Omega\times \R^+)|\ u|_{\partial\Omega}=0\}$, $R$ can be defined as $R(u_0,g)=R_bg+R_0(u_0-R_bg|_{t=0})$ so we focus on the construction of $R_0$.


For $u_0\in H^\theta_0$ $(H^{1/2}_{00}$ for $\theta=1/2$), we extend it as an odd function of $x_d$, $I(u_0)$ defined on $\R^d$. The map $I:\ H^\theta_0(\R^{d-1}\times \R^+)\to H^\theta(\R^d)$ is continuous as it is clearly the case for $\theta=0,1$. Define now
\[
\widehat{R_I(I(u_0))}(\xi,\delta)=\chi(\cxi t)\widehat{I(u_0)}(\xi),
\]
where $\chi$ is as in Lemma~\ref{releve}. According to the proof of Lemma~\ref{releve}, $R_I\circ I:\ H^\theta\to H^{\theta+1/2}(\R^d\times \R^+)$ is continuous, moreover by construction $R_I\circ I(u_0)$ is an odd function of $x_d$, therefore necessarily $R_I\circ I (u_0)|_{x_d=0}=0$. Thus by taking the restriction on $\R^{d-1}\times \R^+_{x_d}
\times \R_t^+$, $R_0:=R_I\circ I$ solves the problem.
\end{proof}


\begin{prop}
For $0<\theta<1$, $[Y_0,Y_1]_\theta=Y_\theta$.
\end{prop}
\begin{proof}
On one hand, the map $(u_0,g)\to u_0(x',x_d)-g(x',x_d)$ is continuous $Y_i\to H^i_0$ for $i=0,1$, therefore by interpolation it is continuous $[Y_0,Y_1]_\theta\to H^\theta_0$. This gives the first inclusion
\begin{equation}\label{inclusion}
[Y_0,Y_1]_\theta\subset Y_\theta.
\end{equation}
On the other hand, from Lions--Peetre reiteration theorem, for any $0<s,\theta<1$
\[
[[Y_0,Y_1]_s,Y_1]_\theta=[Y_0,Y_1]_{\theta+s(1-\theta)}.
\]
If we have for some $s<1/2$, $[Y_s,Y_1]_\theta\supset Y_{\theta+s(1-\theta)}$ for any $0<\theta<1$, then by reiteration this implies $[Y_0,Y_1]_\theta=Y_\theta$ for $\theta >s$. On the other hand, the case $\theta\leq s$ is contained in Lemma~\ref{interpfacile}.

For any $0<r<1$ we define the map
\[
u\in H^{r+1/2}\left(\Omega\times \R_t^+\right)
\to \Tr(u)=\left(u|_{t=0},u|_{x_d=0}\right).
\]
It is easily seen that $\Tr$ is continuous $H^{3/2}\to Y_1$ and $H^{1/2+s}\to Y_s$ for $0<s<1/2$. As it is well known that $[H^{s+1/2},H^{3/2}]_{\theta}=H^{1/2+\theta +(1-\theta)s}$, we deduce by interpolation
\[
\Tr:\ H^{1/2+(1-\theta)s+\theta}=\left[H^{s+1/2},H^{3/2}\right]_\theta\to [Y_s,Y_1]_\theta\text{ is continuous}.
\]
We observe now that the lifting $R$ from Lemma~\ref{relevecoin} is a right inverse for $\Tr$: for fixed $0<s<1/2$ and any $0<\theta<1$, we have $\Tr\circ R=I_d:\ Y_{\theta+s(1-\theta)}
\to Y_{\theta+s(1-\theta)}$. Since $R$ maps $Y_{\theta+s(1-\theta)}$ to $H^{s(1-\theta)+\theta+1/2}$, this implies
\[
Y_{\theta+s(1-\theta)}\subset [Y_s,Y_1]_\theta,
\]
which was the required converse inclusion.
\end{proof}

%\begin{comment}
%\begin{rema}
%We shall use in the sequel the simplified resolvent estimate:
%\[
%\|e^{-\gamma t}u\|_{H^i(\Omega_T)}\lesssim \frac{1}{\gamma}
%\big(\|e^{-\gamma t}f\|_{H^i}+\|e^{-\gamma t}g\|_{H^i} +\|u_0\|_{H^i}\big),\ i=0,1.
%\]
%By interpolation this implies
%\begin{equation}\label{resolvS}
%\|u\|_{H^s(\Omega_T)}\lesssim \frac{1}{\gamma}
%\big(\|(L+\gamma)u\|_{H^s}+\|g\|_{H^s} +\|u_0\|_{H^s}\big),\ 0\leq s\leq 1.
%\end{equation}
%\end{rema}
%\end{comment}

\subsection{The case $s>1$}
We denote $s=k+\theta$, $0\leq \theta<1$. According to the integer case, we already have $u\in \cap C^{k-j}H^j$. For any tangential multi-index $\alpha$ of order $k$ (that is, $\alpha_d=0,|\alpha|=k$), $\partial^\alpha u$ satisfies
\begin{equation}\label{derBVP}
\left\{
\begin{array}{ll}
L(\partial^\alpha u)=\partial^\alpha f+[L,\partial^\alpha]u,\\
B\partial^\alpha u|_{\partial\Omega}=\partial^\alpha g+[B,\partial^\alpha]u|_{\partial\Omega},\\
\partial^\alpha u|_{t=0}=L_\alpha (u_0)+ L_\alpha' (f)|_{t=0}.
\end{array}
\right.
\end{equation}
where $L_\alpha,L_\alpha'$ are differential operators of respective order $\alpha,\alpha-1$. Regularity will again be obtained by regularization of the data, we distinguish three cases:


\subsubsection{The case $0<\theta<1/2$}
With the same argument as in the integer case (note that the condition $\theta<1/2$ allows to use Lemma~\ref{releve}), there exists regularized data $(u_{0,n},g_n,f_n)\in (H^{k+1})^3$, converging to $(u_0,g,f)$ in $(H^k)^3$ that satisfy the compatibility conditions up to order $k+1$. The corresponding solution $u_n$ belongs to $\cap_0^{k+1} C_t^jH^{k+1-j}$ so that we may apply the resolvent estimate~\eqref{resolvreg} to $\partial^\alpha u_n$ with $s=\theta$, combined with basic trace estimates and the commutator estimate $\|[\partial^\alpha,L]u_n\|_{H^\theta_\gamma}\lesssim
\|u_n\|_{H^s_\gamma}$:
\begin{equation}\label{keyresolv}
\gamma\|\partial^\alpha u_n\|_{H^\theta_\gamma}^2 \lesssim \|u_{0,n}\|_{H^s}^2+\|g_n\|_{H^s_\gamma}^2 +\frac{\|f_n\|_{H^s_\gamma}^2+
\|u_n\|_{H^s_\gamma}^2}{\gamma}.
\end{equation}
Due to the boundary being non characteristic, we deduce as for the integer case (note that the fractional regularity gained here includes conormal regularity) for $\gamma$ large enough only depending on $s$, that we have the resolvent estimate
\[
\gamma \|u_n\|_{H^s_\gamma}^2
\lesssim \|u_{0,n}\|^2_{H^s}+\|g_n\|_{H^s_\gamma}^2+
\frac{\|f_n\|_{H^s_\gamma }^2}{\gamma}.
\]
With the resolvent estimate available, the semi group estimate is now an immediate consequence of the case $0<s<1$ applied to~\eqref{derBVP}:
\begin{align*}
\left\|e^{-\gamma t} \partial^\alpha u_n\right\|_{C_tH^\theta}^2&\lesssim
\|u_{0,n}\|_{H^s(\Omega)}^2 +\|f_n\|_{H^s_\gamma([0,T]\times \Omega)}^2 +\|[\partial^\alpha,L]u_n\|_{H^\theta_\gamma(\Omega\times [0,T])}^2 +\|g_n\|_{H^s_\gamma}^2
\\
&\lesssim
\|u_{0,n}\|_{H^s(\Omega)}^2 +\|f_n\|_{H^s_\gamma([0,T]\times \Omega)} +\|g_n\|_{H^s_\gamma}^2.
\end{align*}
Normal regularity is obtained thanks to the boundary being non characteristic as in~\eqref{characgain}. 


Letting $n\to \infty$, we deduce that $e^{-\gamma t}u$ is in $H^s(\R^+\times \Omega)\cap (\cap_{j=0}^kC^j(\R^+,H^{s-j}(\Omega))$ and satisfies the semi group estimate and the resolvent estimate.

\subsubsection{The case $1/2<\theta<1$}
This can be done with exactly the same argument. Actually, the construction of regularized data $(u_{0,n},g_n,f_n)\in (H^{k+1})^3$ that satisfy compatibility conditions up to order $k+1$ and converging to $(u_0,g,f)$ in $(H^s)^3$ is even simpler. Indeed $(u_0,g,f)$ satisfy compatibility conditions up to order $k+1$, hence any regularization of $(u_0,g,f)$ satisfies
\[
\forall\,1\leq j\leq k+1,\
\left\|\underbrace{\partial_t^{j-1}g_n|_{t=0} -\sum_{l=0}^{k-1}\binom{j}{l} \left(\partial_t^lB\right)v_{j-1-l,n}|_{\partial\Omega}}_{:=\varepsilon_{j,n}}\right\|_{H^{s-j+1/2}}\longrightarrow_{n} 0,
\]
and it suffices to modify $g_n$ as $g_n-\delta_n$ where $\delta_n$ is a function in $H^{k+1}(\partial\Omega\times\R_t^+)$ that satisfies for $1\leq j\leq k+1$, $\partial_t^{j-1}\delta_n|_{t=0}=\varepsilon_{j,n}$

\subsubsection{The case $\theta=1/2$}\label{cas1/2}
When $s=k+1/2$, the compatibility conditions are satisfied in particular up to order $k$. From the previous study, we have $e^{-\gamma t}u\in (\cap_{j=0}^k C_t^jH^{k+\theta-j})\cap H^{k+\theta}$ for any $\theta<1/2$, with the estimate
\[
\left\|e^{-\gamma t}u\right\|_{\left(\cap_{j=0}^k C_t^jH^{j+\theta-j}\right)} +\left\|e^{-\gamma t}u\right\|_{H^{k+\theta}}
\leq C(\theta)\|(u_0,g,f)\|_{(H^s)^3}.
\]
Of course this is not enough to conclude, but the estimate can be sharpened: apply estimate~\eqref{resolvsupersharp} to~\eqref{derBVP} for $\theta<1$ and any tangential multi-index $\alpha\in \N^d$, $|\alpha|=k$, this reads
\[
\gamma\left\|\partial^\alpha u\right\|_{H^\theta_\gamma}^2 \lesssim \left\|\left(L_\alpha u_{0}+L_\alpha'f|_{t=0},e^{-\gamma t}\left(\partial^\alpha g+\left[B,\partial^\alpha\right]u|_{\partial\Omega}\right)\right)\right\|_{X_\theta}^2 +\frac{\|f\|_{H^{k+\theta}_\gamma}^2+
\|u\|^2_{H^{k+\theta}_\gamma}}{\gamma}.
\]
Recall that the compatibility conditions at order $j$ are
\[
\forall\,1\leq j\leq k,\
\partial_t^{j-1}g|_{t=0}-\sum_{l=0}^{j-1}
\binom{j}{l}\left(\partial_t^lB\right)v_{j-1-l}|_{\partial\Omega}=0,
\]
and at order $k+1/2$
\begin{equation*}
\partial_t^{k}g(x',t)-\left(\sum_{l=0}^{k}
\binom{k}{l}\left(\partial_t^lB\right)v_{k-1-l}(x',t)\right)
\in H^{1/2}_{00}\left(\R^{d-1}\times (\R^+)\right).
\end{equation*}
\begin{comment}
It is more convenient to rewrite them in term of the weighted function $e^{-\gamma t}u$: define as in~\eqref{taylor} $\clA_\gamma=\gclA-\gamma$, and inductively
\[
v_{0,\gamma}=u_0,\ v_{j,\gamma}=\sum_{l=0}^j\binom{j}{l}(\partial_t^l\clA_\gamma|_{t=0}) v_{j-l}+\partial_t^j(e^{-\gamma t}f)|_{t=0}.
\]
The compatibility conditions rewrite
\begin{equation}\label{CCgamma}
\forall\,1\leq j\leq k,\
\partial_t^{j-1}(e^{-\gamma t}g)|_{t=0}-\sum_{l=0}^{j-1}
\binom{j}{l}(\partial_t^lB)v_{j-1-l,\gamma}|_{\partial\Omega}=0,
\end{equation}
and
\begin{equation}\label{CGgamma}
\partial_t^{k}(e^{-\gamma t}g(x',t))-\left(\sum_{l=0}^{k}
\binom{k}{l}(\partial_t^lB)v_{k-1-l,\gamma}(x',t)\right)
\in H^{1/2}_{00}\left(\R^{d-1}\times (\R^+)\right).
\end{equation}
\end{comment}

As a consequence, for any $j\leq k+1$ and any $\beta\in \N^{d-1},\ |\beta|=k+1-j$,
\begin{equation}\label{compafrac}
\partial_{x'}^\beta\partial_t^{j-1}g(x',t) -\partial_{x'}^\beta\sum_{l=0}^{j-1}
\binom{j-1}{l}\left(\partial_t^lB\right)v_{j-1-l}(x',t)\in H^{1/2}_{00}\left(\R^{d-1}\times \R^+\right).
\end{equation}
Furthermore, $e^{-\gamma t}u\in H^k(\Omega
\times \R_t^+)$, hence for any multi-index of order $k-1$
\begin{equation}\label{compagratuite}
\left\|e^{-\gamma t}\partial^\alpha u|_{x_d=0} -e^{-\gamma t}\partial^\alpha u|_{t=0}\right\|_{H^{1/2}_{00}\left(\R^{d-1}\times \R^+\right)}
\lesssim \left\|e^{-\gamma t}u\right\|_{H^{k}\left(\R^{d-1}\times (\R^+)^2\right)}.
\end{equation}
Now to make~\eqref{derBVP} more explicit, let us write $\partial^\alpha=
\partial_t^{j}\partial_{x'}^\beta$, $\beta\in
\N^{d-1},\ |\beta|=k-j$. Then $\partial^\alpha u|_{t=0} =\partial_{x'}^\beta v_{j}\in H^{1/2}(\R^{d-1}\times \R^+)$, the compatibility condition of order $1/2$ for~\eqref{derBVP} is thus
\[
e^{-\gamma t}\left(\partial^\alpha g+ [B,\partial^\alpha]u|_{\partial\Omega}\right) -B\partial_{x'}^\beta v_{j}\in H^{1/2}_{00}\left(\R^{d-1}\times \R^+\right).
\]
With basic computations, we now check that it is implied by~\eqref{compafrac},\eqref{compagratuite}:
\begingroup
\allowdisplaybreaks
\begin{multline*}
e^{-\gamma t}\left(\partial^\alpha g + \left[B,\partial^\alpha\right]u|_{\partial\Omega}\right)-B\partial_{x'}^\beta v_{j}
\\
\mkern-10mu{\begin{aligned}
&= e^{-\gamma t}\left(\partial^\alpha g-\partial_{x'}^\beta
\sum_{l=0}^j\binom{j}{l}\left(\partial_t^lB\right)\partial_t^{j-l}u|_{\partial\Omega}\right) +e^{-\gamma t}B\partial^\alpha u|_{\partial\Omega}-B\partial_{x'}^\beta v_{j}
\\
&=e^{-\gamma t}\partial^\alpha g -\partial_{x'}^\beta \sum_1^j\binom{j}{l}\!\!\left(\partial_t^lB\right) v_{j-l} -B\partial_{x'}^\beta v_{j}
-\partial_{x'}^\beta\sum_{l=1}^j
\binom{j}{l}\!\!\left(\partial_t^lB\right)
\!\!\left(e^{-\gamma t}\partial_t^{j-l}u|_{\partial\Omega} -v_{j-l}\right)\\
& -\partial_{x'}^\beta \left(Be^{-\gamma t}\partial_t^j u|_{\partial\Omega}\right) +e^{-\gamma t}B\partial^\alpha u|_{\partial\Omega}\\
&=e^{-\gamma t}\left(\partial^\alpha g -\partial_{x'}^\beta \sum_{l=0}^j\binom{j}{l}\left(\partial_t^lB\right)v_{j-l}\right)
 -\partial_{x'}^\beta\sum_{l=1}^j
\binom{j}{l}\left(\partial_t^lB\right)
\left(e^{-\gamma t}\partial_t^{j-l}u|_{\partial\Omega}-v_{j-l}\right)\\
&+\left[B,\partial_{x'}^\beta\right]e^{-\gamma t}\partial_t^ju|_{\partial\Omega} -\left[B,\partial_{x'}^\beta\right]v_{j}.
\end{aligned}}
\end{multline*}
\endgroup
For $j\leq k$, due to the compatibility condition~\eqref{compafrac}, in the last equality, the first term in the first line is in $H^{1/2}_{00}$. The $H^{1/2}_{00}$ norm of the second term is easily controlled by writing
\[
e^{-\gamma t}\partial_t^{j-l}u|_{\partial\Omega}-v_{j-l} =e^{-\gamma t}\left(\partial_t^{j-l}u|_{\partial\Omega}- v_{j-l}\right)+\left(1-e^{-\gamma t}\right)v_{j-l},
\]
the first term can be bounded thanks to~\eqref{compagratuite} while for the second one we simply use $(1-e^{-\gamma t})/t
\lesssim 1$. The same argument is used for the second line. We deduce that for $\theta <1/2$, $\alpha$ tangential, $|\alpha|\leq k$
\begin{multline*}
\gamma\left\|\partial^\alpha u\right\|^2 _{H^{\theta}_\gamma\left(\Omega\times \R_t^+\right)} \\
\lesssim C(\gamma)\left(\|(u_0,g,f)\|_{\left(H^{k+1/2}\right)^3} +\left\|g -\sum_0^k\binom{k}{l}\left(\partial_t^lB\right) v_{k-1-l}\right\|_{H^{1/2}_{00} \left(\R^{d-1}\times \R^+\right)}\right)\\
+\frac{\|u\|_{H^{k+\theta}_\gamma}^2}{\gamma}.
\end{multline*}
Using that the boundary is non characteristic, we recover
\[
\gamma\|u\|^2 _{H^{k+\theta}_\gamma\left(\Omega\times \R_t^+\right)} \lesssim \|(u_0,g,f)\|_{\left(H^{k+1/2}\right)^3} +\left\|g -\sum_0^k\binom{k}{l}(\partial_t^lB) v_{k-1-l}\right\|_{H^{1/2}_{00} \left(\R^{d-1}\times \R^+\right)}.
\]
This estimate is uniform in $\theta <1/2$, we deduce that the same estimate holds for $\theta =1/2$. Finally we deduce that the semi group estimate is true with the same argument as for the end of the case $0<\theta <1/2$ : consider the problem~\eqref{derBVP}, since the commutator $[L,\partial^\alpha]u$ belongs to $H^{1/2}$, the semi group estimate of the case $s=1/2$ can be applied to bound $\|\partial^\alpha u\|_{C_tH^{1/2}}$, $\alpha$ any tangential multi-index of order $k$, the normal regularity follows from the usual argument.


This ends the proof of Theorem~\ref{mainth}.



\begin{comment}
Oui mais non, le forçage $\partial_t f+[L,\partial_t]u$ n'est a priori que $L^2$ et ça coince.


Un remède ? Pour $\theta<1/2$, étant donné des données $(u_0,g,f)$ satisfaisant la relation de compa à l'ordre $1$, on produit $(u_{0,n},f_n,g_n)$ dans $H^2$ avec $\n u_{0,n}|_{\partial\Omega}=0$, $f_n|_{x=t=0}=0$, $\partial_tg|_{t=0}=0$ et convergeant vers les données: il suffit de reprendre le cas entier : on prend $(u_{0,n},f_n,g_n)$ dans $H^2$ sans les conditions de compatibilité, on a
\[
\|Bu_{0,n}-g_n\|_{H^{1/2+\theta}}<<1,\ \partial_tBu_{0,n} +B(Au_{0,n}+f_n)|_{\partial\Omega}-\partial_tg_n|_{t=0}
\in H^{1/2} \text{ (non petit)}.
\]
On modifie $g_n$ par $\widetilde{g_n}=g_n+R_0(Bu_{0,n}-g_n) +R_1(\partial_tBu_{0,n} +B(Au_{0,n}+f_n)|_{\partial\Omega}-\partial_tg_n$, $R_1$ choisi de sorte que la norme $H^{1+\theta}$ soit petite. En oubliant le $n$, on a
\[
\|\partial_t u\|_{L^\infty_tH^\theta}\lesssim \|\partial_t f\|_{H^\theta} +\|[L,\partial_t]u\|_{H^\theta}+\|Au_0+f|_{t=0}\|_{H^\theta} +\|\partial_t g\|_{H^\theta}.
\]
Bon mais là pas bien, en ajoutant les autres dérivées tangentielles, on contrôle une régularité $C^1H^{\theta}\cap C^0H^{1+\theta}$ par la régularité $H^{1+\theta}(\Omega_T)$, pas suffisant. À la place on peut commencer par l'estimée de résolvante~\ref{resolvS}
\[
\|\partial_t(e^{-\gamma t}u)\|_{H^s}\lesssim
\frac{1}{\gamma}\big(\|\partial_t e^{-\gamma t}f\|_{H^\theta} +\|[L,\partial_t]e^{-\gamma t}u\|_{H^\theta}+\|Au_0+f|_{t=0}\|_{H^\theta} +\|\partial_t e^{-\gamma t}g\|_{H^\theta}\big).
\]
La même estimation tangentielle donne pour $1\leq j\leq d-1$
\[
\|\partial_j(e^{-\gamma t}u)\|_{H^s}\lesssim
\frac{1}{\gamma}\big(\|\partial_j e^{-\gamma t}f\|_{H^\theta} +\|[L,\partial_t]e^{-\gamma t}u\|_{H^\theta}+\|Au_0+f|_{t=0}\|_{H^\theta} +\|\partial_j e^{-\gamma t}g\|_{H^\theta}\big).
\]
Finalement, en utilisant la non caractéristicité
\[
\|\partial_d e^{-\gamma t}u\|_{H^\theta}\lesssim \|(\partial_t+\sum_1^{d-1}e^{-\gamma t}u\|_{H^\theta}\lesssim \frac{1}{\gamma}
\|e^{-\gamma t}u\|_{H^{\theta +1}}+
\|u_0\|_{H^{\theta+1}}+\|g\|_{H^{\theta+1}}
\|f\|_{H^{\theta+1}},
\]
on peut maintenant manger $\frac{1}{\gamma}\|e^{-\gamma t}u\|_{H^{\theta +1}}$ et obtenir la bonne estimation
\[
\|u\|_{H^{\theta+1}(\Omega_T)}\lesssim \|g\|_{H^{\theta+1}} +\|f\|_{H^{\theta +1}}+\|u_0\|_{H^{\theta+1}}.
\]
Ce gain de régularité suffit pour la régularité $C_tH^{1+\theta}$ puisque maintenant $[\partial,L]u\in H^\theta(\Omega_T)$.

On se souvient que $u$ est en fait $u_n$, faisant $n\to \infty$, on conclut
\begin{prop}
Si $Bu_0|_{\partial\Omega}=g$, $1<s<3/2$ La solution de
\begin{equation}
\left\{
\begin{array}{ll}
Lu=f,\\
u|_{t=0}=u_0,\\
Bu|_{\partial\Omega_T}=g.
\end{array}
\right.
\end{equation}
la solution est $H^s(\Omega_T)\cap C_TH^s$ (et même si on veut de la régularité h\"oldérienne en temps mais on s'en fiche).
\end{prop}

Pour $3/2<s<2$, on peut faire le même genre d'approximation par relèvement, on s'en passera. Pour $s=3/2$ c'est délicat, car il n'est pas clair qu'on puisse approcher $u_0,g,f$ par des données satisfaisant la condition de compatibilité d'ordre supérieur et respectant la topologie bizarre type Lions--Magenes.


\`A défaut, on peut faire cet argument : si les données satisfont les CC d'ordre $3/2$, on a $u\in H^1(\Omega_T)$ et on peut écrire
\begin{align*}
\|\partial_t e^{-\gamma t}u\|_{H^{\theta}(\Omega_T)}&\lesssim
\frac{\|\partial_t e^{-\gamma t}f\|_{H^{\theta}}}{\gamma}+
\|Au_0+f|_{t=0},\partial_t g-(\partial_t B)u\|_{[X_0,X_1]_\theta}
\\
&\lesssim \frac{\|\partial e^{-\gamma t}f\|_{H^{\theta}}}{\gamma}+
\|Au_0+f|_{t=0},\partial_t g-(\partial_t B)u\|_{[X_0,X_1]_{1/2}}
\end{align*}
Comme avant, on en tire $\|u\|_{H^{1+\theta}}\lesssim
\|u_0\|_{H^{3/2}}+\|f\|_{H^{3/2}} +\|\partial_t (Bu),Au_0+f\|_{[X_0,X_1]_{1/2}}$. L'estimée est clairement indépendante de $\theta <1/2$, on en tire par passage à la limite l'estimation pour $\theta=1/2$. L'estimation ponctuelle en temps s'obtient comme dans le cas précédent.
\end{comment}



\begin{thebibliography}{99}
\bibitem{Amann2}Amann, Herbert Linear and quasilinear parabolic problems. Vol. II: Function spaces, Monographs in Mathematics, 106, Birkhäuser; Springer, 2019 | MR | Zbl | DOI
\bibitem{Audiard1}Audiard, Corentin On mixed initial-boundary value problems for systems that are not strictly hyperbolic, Appl. Math. Lett., Volume 24 (2011) no. 5, pp. 757-761 | DOI | MR | Zbl
\bibitem{Audiard7}Audiard, Corentin Global Strichartz estimates for the Schrödinger equation with non zero boundary conditions and applications, Ann. Inst. Fourier, Volume 69 (2019) no. 1, pp. 31-80 | MR | Zbl | DOI | Numdam
\bibitem{Benzoni3}Benzoni-Gavage, Sylvie; Serre, Denis Multidimensional hyperbolic partial differential equations. First-order systems and applications, Oxford Mathematical Monographs, The Clarendon Press; Oxford University Press, 2007 | MR | Zbl
\bibitem{ChaPi}Chazarain, Jacques; Piriou, Alain Introduction à la théorie des équations aux dérivées partielles linéaires, Gauthier-Villars, 1981 | MR | Zbl
\bibitem{Friedrichs}Friedrichs, Kurt O. Symmetric positive linear differential equations, Commun. Pure Appl. Math., Volume 11 (1958), pp. 333-418 | DOI | MR | Zbl
\bibitem{Grisvard}Grisvard, Pierre Caractérisation de quelques espaces d’interpolation, Arch. Ration. Mech. Anal., Volume 25 (1967), pp. 40-63 | DOI | MR | Zbl
\bibitem{hormanderBVP}Hörmander, Lars Pseudo-differential operators and non-elliptic boundary problems, Ann. Math., Volume 83 (1966), pp. 129-209 | DOI | MR | Zbl
\bibitem{lanigu}Iguchi, Tatsuo; Lannes, David Hyperbolic free boundary problems and applications to wave-structure interactions, Indiana Univ. Math. J., Volume 70 (2021) no. 1, pp. 353-464 | DOI | MR | Zbl
\bibitem{Kreiss}Kreiss, Heinz-Otto Initial boundary value problems for hyperbolic systems, Commun. Pure Appl. Math., Volume 23 (1970), pp. 277-298 | MR | Zbl | DOI
\bibitem{lionsmagenes}Lions, Jacques-Louis; Magenes, Enrico Problèmes aux limites non homogènes et applications. Vol. 1, Travaux et Recherches Mathématiques, 17, Dunod, 1968 | MR | Zbl
\bibitem{lionsmagenes2}Lions, Jacques-Louis; Magenes, Enrico Problèmes aux limites non homogènes et applications. Vol. 2, Travaux et Recherches Mathématiques, 18, Dunod, 1968 | MR | Zbl
\bibitem{LaxPhillips}Lax, Peter D.; Phillips, Ralph S. Local boundary conditions for dissipative symmetric linear differential operators, Commun. Pure Appl. Math., Volume 13 (1960), pp. 427-455 | DOI | MR | Zbl
\bibitem{Met}Métivier, Guy The block structure condition for symmetric hyperbolic systems, Bull. Lond. Math. Soc., Volume 32 (2000) no. 6, pp. 689-702 | DOI | MR | Zbl
\bibitem{Met2}Métivier, Guy Stability of multidimensional shocks, Advances in the theory of shock waves (Progress in Nonlinear Differential Equations and their Applications), Volume 47, Birkhäuser, 2001, pp. 25-103 | MR | Zbl | DOI
\bibitem{Metsemigroupe}Métivier, Guy On the ${L}^2$ well posedness of hyperbolic initial boundary value problems, Ann. Inst. Fourier, Volume 67 (2017) no. 5, pp. 1809-1863 | MR | Zbl | DOI | Numdam
\bibitem{Rauch}Rauch, Jeffrey B. ${\mathcal{L}}_{2}$ is a continuable initial condition for Kreiss’ mixed problems, Commun. Pure Appl. Math., Volume 25 (1972), pp. 265-285 | MR | Zbl | DOI
\bibitem{RauchMassey}Rauch, Jeffrey B.; Massey, Frank J. III Differentiability of solutions to hyperbolic initial-boundary value problems, Trans. Am. Math. Soc., Volume 189 (1974), pp. 303-318 | MR | Zbl
\bibitem{Tartakoff}Tartakoff, David S. Regularity of solutions to boundary value problems for first order systems, Indiana Univ. Math. J., Volume 21 (1972), pp. 1113-1129 | DOI | MR | Zbl
\end{thebibliography}
\end{document}
